

In this experiment, comprehensive Quality Assurance (QA) testing of a diagnostic computed tomography (CT) machine (GE Optima CT580W) was performed in accordance with regulatory guidelines to verify its mechanical accuracy, dosimetric output, image quality, and radiation safety. The evaluations were conducted utilizing specialized equipment: the Pro-CT MINI phantom, RTI Piranha multi-meter, and RTI CT Ionization Chamber. 

The assessment of mechanical and geometric fidelity revealed that the gantry tilt indicator is highly accurate, showing a maximum deviation of only $0.8^\circ$ (within the acceptable $\pm 2^\circ$ tolerance). The slice thickness verification corroborated these findings; for a nominal $5\text{ mm}$ slice, the average measured thickness was $4.325\text{ mm}$, securely fulfilling the AERB tolerance requirement of $\pm 1\text{ mm}$.

Output performance parameters—specifically the accuracy of the operating potential (kVp), timer accuracy, output linearity, and output reproducibility—demonstrated consistent and stable generator performance throughout the clinical range. Furthermore, the dosimetric assessment measuring the Weighted CT Dose Index ($CTDI_w$) confirmed that patient radiation exposures align predictably with the manufacturer's operational specifications and normative tolerance bounds.

Image quality was strictly evaluated using the Pro-CT MINI phantom:
\begin{itemize}
    \item \textbf{Low Contrast Resolution:} The system successfully resolved low-contrast details of $2\text{ mm}$ at a $0.5\%$ contrast difference, surpassing the $2.5\text{ mm}$ minimum expected tolerance.
    \item \textbf{High Contrast Resolution:} Spatial detail verification resolved line pair patterns up to $5\text{ lp/cm}$, comprehensively satisfying the minimum clinical requirement ($3.12\text{ lp/cm}$).
\end{itemize}

Finally, radiation leakage mapping and the radiation protection survey around the installation showed conclusively that stray and transmitted radiation levels around the scanner boundaries and housing are safely restricted within established regulatory limits.

Overall, all evaluated physical and functional parameters of the diagnostic CT scanner fell safely within operational tolerances. The equipment functions reliably, efficiently maintains patient radiation safety protocols, and continues to produce images of high diagnostic fidelity.